\documentclass[paper={a5paper},maintitle={Layer Layout Test}]{ulsbs-songbook}
\begin{document}
\begin{songs}{}
\beginsong{First Layer Song}[id={fixture-1},key={C}]
  \beginverse
    \[Am]Opening words \[E7]with longer lyric phrases \[F#m/C#]and slash chords.
    \[<* C D;E F @ .3em>Asus2]Combined melody \[<A _ C;E _ G>]empty slots.
    \[<!C>]Zero advance \[</>Am]unpitched \[<* @ 1em>]shared beat.
    \[<C D @ -.4em>]Negative offset \[G]\[<*>]adjacent annotations.
  \endverse
  {\melodyStyleMapping{swapped}\melodySecondaryPosition{below}
    \beginverse
      \[<C# D&;A C>Am7]Swapped secondary notes and accidental hints.
    \endverse}
  \notesoff
  \beginverse
    \[<* C>G]Hidden notes with a visible beat.
  \endverse
  \noteson
  \beginverse
    \[C\kern.5em]Dimension argument \[C\textsuperscript7]unbraced suffix.
    \[F\shrp{}m/C\shrp{}]Explicit accidental groups \[F#\textdegree7]degree suffix.
  \endverse
  \newchords{layeradjacent}
  \beginverse+\memorize[layeradjacent]
    \[Am]\[<*>]Stored adjacent chord and beat \[<C>]\[<*>]\[<D>]\[G]melody.
  \endverse
  \beginverse+\replay[layeradjacent]
    \[<E>^]\[<*>^]\[<F>^]\[^]Explicit adjacent replay annotations.
  \endverse
  \musicnote{Prose \ac{1}{C} example outside verse notation.}
\endsong
\beginsong{Second Layer Song}[id={fixture-2},key={G}]
  \beginverse
    \[G]Independent song \[<G B;D F#>D7]with independent melody.
  \endverse
\endsong
\end{songs}
\end{document}
